\subsection{Optionale Capabilities}
\label{subsec:optionale-capabilities}

Der Katalog enthält sechs Features. \texttt{tauri} benötigt
\texttt{frontend}; \texttt{cloud} und \texttt{database} benötigen
\texttt{backend}; \texttt{postgres} benötigt \texttt{database}. Die
gerichteten Regeln in Abbildung~\ref{fig:feature-dependencies} beschreiben
keine Laufzeitkommunikation
\parencite{kleiveist2026profiles,kleiveist2026database}.

\begin{figure}[H]
  \centering
  \begin{tikzpicture}
    \node[caseBox,minimum width=31mm] (backend) at (0,0) {\texttt{backend}};
    \node[caseBox,minimum width=31mm] (frontend) at (7,0) {\texttt{frontend}};

    \node[caseBox,minimum width=31mm] (cloud) at (-2,-1.6) {\texttt{cloud}};
    \node[caseOptional,minimum width=31mm] (database) at (2,-1.6)
      {\texttt{database}\\optional};
    \node[caseBox,minimum width=31mm] (tauri) at (7,-1.6) {\texttt{tauri}};
    \node[caseOptional,minimum width=31mm] (postgres) at (2,-3.2)
      {\texttt{postgres}\\auswählbar};

    \draw[caseFlow] (cloud) -- (backend);
    \draw[caseFlow] (database) -- (backend);
    \draw[caseFlow] (postgres) -- (database);
    \draw[caseFlow] (tauri) -- (frontend);
    \node[font=\sffamily\scriptsize,text=black!60,below=3mm of postgres]
      {Pfeilrichtung: benötigt};
  \end{tikzpicture}
  \caption{Implementierte Abhängigkeiten des Feature-Katalogs}
  \label{fig:feature-dependencies}
\end{figure}


Nur \texttt{postgres} ist derzeit als auswählbare Capability markiert. Für
Backendprofile löst sie zunächst das optionale, intern bleibende
\texttt{database} auf und ergänzt SQLAlchemy, Alembic, Psycopg sowie Tests. Der
PostgreSQL-Dienst ist im Compose-Dokument bereits als inaktives Profil
deklariert; die Capability ergänzt Backend-Infrastruktur und Abhängigkeiten
\parencite{kleiveist2026deployment,kleiveist2026governance-snapshot}.
So besteht \texttt{web-cloud + postgres} aus Frontend, Backend, Cloud,
Datenbankbasis und PostgreSQL-Erweiterung. Die Auswahl aktiviert also nicht
erst den Compose-Eintrag, sondern die zugehörigen Backend-Pfade.

Der Resolver ergänzt nur optionale Voraussetzungen. Deshalb werden
\texttt{web-only + postgres} und \texttt{desktop-local + postgres} wegen des
fehlenden Backends abgewiesen. Neue Katalogeinträge gelten nicht allein durch
ihre Deklaration als implementiert.
Sie benötigen zusätzlich eigene Pfade, Abhängigkeiten und Tests.

Zwei Implementierungsabweichungen bleiben: Der Dialog bietet nur
\texttt{selectable} PostgreSQL an, eine direkte \texttt{--with database}-Angabe
akzeptiert jedoch die für Backendprofile optionale Datenbankfähigkeit. Zudem
stammen Einträge für \path{.env.example} entgegen der allgemeinen
Profildokumentation aus \path{config/environment.toml}, nicht aus
\path{features.toml}. Die Presets bleiben unverändert; die Trennung zwischen
interner und auswählbarer Capability ist aber nicht vollständig durchgesetzt.
\beginnerinput{workspace/statement/400_projektprofile/470}
